\documentclass[12pt]{article}

\usepackage[utf8]{inputenc}
\usepackage{amsmath, amssymb, amsthm}
\usepackage{geometry}
\geometry{a4paper, margin=2.5cm}

\title{Analiză Funcțională — Exercițiu}
\author{}
\date{}

\theoremstyle{definition}
\newtheorem{definition}{Definiție}

\theoremstyle{plain}
\newtheorem{proposition}{Propoziție}
\newtheorem{theorem}{Teoremă}

\begin{document}

\maketitle


\begin{definition}
Fie


\[
\ell^1 = \left\{ x = (x_n)_{n\ge 1} : \sum_{n=1}^\infty |x_n| < \infty \right\},
\qquad
\|x\|_1 = \sum_{n=1}^\infty |x_n|.
\]



Definim


\[
Y = \left\{ x \in \ell^1 : \exists N \in \mathbb{N} \text{ astfel încât } x_n = 0 \text{ pentru toate } n \ge N \right\}.
\]


\end{definition}

\begin{proposition}
Subspațiul $Y$ nu este închis în $\ell^1$.
\end{proposition}

\begin{proof}
Pentru orice $x = (x_n) \in \ell^1$ definim șirurile trunchiate


\[
x^{(N)} = (x_1, x_2, \dots, x_N, 0, 0, \dots).
\]


Observăm că $x^{(N)} \in Y$ pentru orice $N$.

Calculăm norma diferenței:


\[
\|x - x^{(N)}\|_1 = \sum_{n>N} |x_n|.
\]



Deoarece $x \in \ell^1$, restul seriei tinde la zero:


\[
\sum_{n>N} |x_n| \longrightarrow 0 \quad \text{când } N \to \infty.
\]



Prin urmare,


\[
x^{(N)} \to x \quad \text{în } \ell^1.
\]



Aceasta arată că pentru orice element din $\ell^1$, există un șir convergent din $Y$ care să tindă la el, deci


\[
\overline{Y} = \ell^1.
\]



Dar știm că $Y \ne \ell^1$, deci $Y$ nu poate fi închis.
\end{proof}

\begin{proposition}
Pentru fiecare $n \in \mathbb{N}^\ast$, definim aplicația liniară


\[
T_n : Y \to \mathbb{R}, \qquad T_n(u) = n\,u_n, \quad u = (u_k)_{k\ge 1} \in Y.
\]


Atunci $T_n \in Y^*$, iar norma lui $T_n$ este


\[
\|T_n\|_{Y^*} = n.
\]


\end{proposition}

\begin{proof}
Mai întâi observăm că $T_n$ este liniară, deoarece pentru orice $u,v \in Y$ și $\alpha,\beta \in \mathbb{R}$ avem


\[
T_n(\alpha u + \beta v)
= n(\alpha u_n + \beta v_n)
= \alpha n u_n + \beta n v_n
= \alpha T_n(u) + \beta T_n(v).
\]



Arătăm acum că $T_n$ este continuă pe $(Y,\|\cdot\|_1)$. Pentru orice $u \in Y$,


\[
|T_n(u)|
= |n u_n|
\le n \sum_{k=1}^\infty |u_k|
= n \|u\|_1.
\]


Rezultă că


\[
\|T_n\|_{Y^*}
= \sup_{\|u\|_1 = 1} |T_n(u)|
\le n.
\]



Pentru a obține egalitatea, considerăm vectorul


\[
u^{(n)} = (0,\dots,0,\underbrace{1}_{\text{poziția } n},0,0,\dots),
\]


care aparține lui $Y$ și satisface $\|u^{(n)}\|_1 = 1$. Atunci


\[
T_n(u^{(n)}) = n \cdot 1 = n,
\]


deci


\[
\|T_n\|_{Y^*}
\ge |T_n(u^{(n)})|
= n.
\]



Combinând cele două inegalități, obținem


\[
\|T_n\|_{Y^*} = n.
\]


Prin urmare, $T_n$ este liniară și continuă pe $Y$, adică $T_n \in Y^*$, și are norma \\$\|T_n\|_{Y^*} = n$.


\end{proof}

\begin{proposition}
Pentru orice $v \in Y$, șirul $\bigl(\lvert T_n(v)\rvert\bigr)_{n \in \mathbb{N}^\ast}$ este mărginit.
\end{proposition}

\begin{proof}
Fie $v = (v_k)_{k\ge 1} \in Y$. Prin definiția lui $Y$, există $N \in \mathbb{N}^\ast$ astfel încât


\[
v_k = 0 \quad \text{pentru toate } k \ge N.
\]


Pentru fiecare $n \in \mathbb{N}^\ast$, avem


\[
T_n(v) = n\,v_n,
\]


deci


\[
|T_n(v)| = |n v_n|.
\]



Dacă $n \ge N$, atunci $v_n = 0$ și, prin urmare,


\[
|T_n(v)| = |n v_n| = 0.
\]


Astfel, șirul $\bigl(|T_n(v)|\bigr)_{n\in\mathbb{N}^\ast}$ are doar un număr finit de termeni nenuli, și anume pentru $n \in \{1,2,\dots,N-1\}$.

Definim


\[
M := \max_{1 \le n \le N-1} |T_n(v)| = \max_{1 \le n \le N-1} |n v_n|.
\]


Acest maxim există deoarece mulțimea $\{1,\dots,N-1\}$ este finită. Atunci, pentru orice $n \in \mathbb{N}^\ast$,


\[
|T_n(v)| \le M.
\]



Rezultă că șirul $\bigl(|T_n(v)|\bigr)_{n\in\mathbb{N}^\ast}$ este mărginit.
\end{proof}

\begin{proposition}
Nu putem aplica teorema Banach--Steinhaus pentru a deduce că șirul


\[
\bigl(\|T_n\|_{Y^*}\bigr)_{n\in\mathbb{N}^\ast}
\]


este mărginit.
\end{proposition}

\begin{proof}
Teorema Banach--Steinhaus (principiul borelienelor uniforme) se aplică unei familii de operatori liniari continui


\[
\mathcal{F} \subset \mathcal{L}(X,Z),
\]


unde $X$ este un spațiu Banach, iar $Z$ este un spațiu normat. Ipoteza este că pentru fiecare $x \in X$, mulțimea


\[
\{ \|T(x)\| : T \in \mathcal{F} \}
\]


este mărginită; concluzia este că


\[
\sup_{T \in \mathcal{F}} \|T\| < \infty.
\]



În situația noastră, familia este $(T_n)_{n\in\mathbb{N}^\ast}$, cu


\[
T_n \in Y^*, \qquad T_n : Y \to \mathbb{R},
\]


iar spațiul de plecare este $Y$, înzestrat cu norma restricție a lui $\|\cdot\|_1$.

Observăm însă că $Y$ \emph{nu este Banach}: în punctul a) am arătat că


\[
\overline{Y}^{\,\ell^1} = \ell^1 \neq Y,
\]


deci $Y$ nu este complet. Prin urmare, ipoteza esențială a teoremei Banach--Steinhaus (completitudinea spațiului de plecare) nu este satisfăcută, și teorema nu este aplicabilă în această formă.

Mai mult, am calculat deja că


\[
\|T_n\|_{Y^*} = n \quad \text{pentru fiecare } n \in \mathbb{N}^\ast,
\]


deci șirul $(\|T_n\|_{Y^*})_{n\in\mathbb{N}^\ast}$ este nemărginit. Dacă am putea aplica teorema Banach--Steinhaus, am obține că acest șir este mărginit, ceea ce ar contrazice calculul direct al normelor. Această contradicție confirmă încă o dată că ipotezele teoremei (în special completitudinea lui $Y$) nu sunt îndeplinite.
\end{proof}

\begin{proposition}
Afirmația de la punctul c) nu mai este valabilă dacă înlocuim $Y$ cu $\ell^1$. Mai precis, există $v \in \ell^1$ pentru care șirul


\[
\bigl(|T_n(v)|\bigr)_{n\in\mathbb{N}^\ast}
\]


nu este mărginit, unde $T_n : \ell^1 \to \mathbb{R}$ este definit prin $T_n(u) = n\,u_n$.
\end{proposition}

\begin{proof}
Definim $v = (v_k)_{k\ge 1}$ prin





\[
v_k =
\begin{cases}
\dfrac{1}{m^2}, & \text{dacă } k = m! \text{ pentru un } m \in \mathbb{N}^\ast,  \\[20pt]
0, & \text{altfel}.
\end{cases}
\]






Atunci


\[
\sum_{k=1}^\infty |v_k|
= \sum_{m=1}^\infty \frac{1}{m^2} < \infty,
\]


deci $v \in \ell^1$.

Pentru $n = m!$ avem


\[
T_{m!}(v) = m! \, v_{m!} = m! \cdot \frac{1}{m^2},
\]


de unde


\[
|T_{m!}(v)| = \frac{m!}{m^2} \xrightarrow[m\to\infty]{} \infty.
\]



Prin urmare, șirul $\bigl(|T_n(v)|\bigr)_{n\in\mathbb{N}^\ast}$ nu este mărginit, ceea ce arată că proprietatea de la punctul c) nu se extinde la întregul spațiu $\ell^1$.
\end{proof}


\end{document}
